This section provides step-by-step instructions for setting up Sanguosuo for local development. The instructions assume a Linux or macOS host; Windows users should substitute paths appropriately.

\subsection{Prerequisites}

Before starting, ensure the following tools are installed:

\begin{itemize}
  \item \textbf{Docker} and \textbf{Docker Compose} (for the backend and database)
  \item \textbf{Android Studio} (recommended) or a standalone Android SDK installation
  \item \textbf{Git} (to clone the repository)
  \item Optionally, \textbf{MikTeX} and \textbf{Perl} with \texttt{latexmk} for compiling the report
\end{itemize}

\subsection{Cloning the Repository}

\begin{verbatim}
git clone https://github.com/Quach-Thien-Lac/cs426-seminar.git
cd cs426-seminar
\end{verbatim}

\subsection{Environment Variables}

Copy the example environment file and fill in the required values:

\begin{verbatim}
cp .env.example .env
\end{verbatim}

Open \texttt{.env} in a text editor. The file requires the following variables:

\begin{center}
\adjustbox{max width=\linewidth}{%
\begin{tabular}{lp{8cm}}
\toprule
\textbf{Variable} & \textbf{Description} \\
\midrule
\texttt{LOCAL\_DB\_PASSWORD}   & Root password for the local MySQL container \\
\texttt{DB\_HOST}              & Database hostname (use \texttt{db} when running inside Docker Compose) \\
\texttt{DB\_PORT}              & MySQL port (default: \texttt{3306}) \\
\texttt{DB\_NAME}              & Database name (use \texttt{Sanguosuo}) \\
\texttt{DB\_USER}              & MySQL username \\
\texttt{DB\_PASSWORD}          & MySQL password for the application user \\
\texttt{SERVER\_PORT}          & Port the Express server listens on (default: \texttt{8080}) \\
\texttt{JWT\_SECRET}           & Secret used to sign JWT tokens \\
\bottomrule
\end{tabular}%
}
\end{center}

\subsection{Backend and Database}

\subsubsection{Starting the Containers}

From the repository root:

\begin{verbatim}
docker compose up
\end{verbatim}

This builds the Express server image from \texttt{server/Dockerfile} and pulls the \texttt{mysql:oraclelinux9} image. The API server will be available at \texttt{http://localhost:8080}. To run the containers in the background (detached mode):

\begin{verbatim}
docker compose up -d
\end{verbatim}

\subsubsection{Seeding the Database}

The \texttt{db/} directory contains two SQL files: \texttt{schema.sql} (table definitions) and \texttt{seed.sql} (sample hero, skill, and reference data). After the containers are running, seed the database as follows:

\begin{enumerate}
  \item Enter the database container shell:
  \begin{verbatim}
  docker compose exec db sh
  \end{verbatim}
  \item Connect to MySQL:
  \begin{verbatim}
  mysql -h 127.0.0.1 -u root -p
  \end{verbatim}
  Enter the value you set for \texttt{LOCAL\_DB\_PASSWORD} when prompted.
  \item Source the schema file (creates all tables):
  \begin{verbatim}
  SOURCE /sql/schema.sql;
  \end{verbatim}
  \item Source the seed file (inserts 148 heroes, 268 skills, and related data):
  \begin{verbatim}
  SOURCE /sql/seed.sql;
  \end{verbatim}
  \item Exit MySQL and the container shell:
  \begin{verbatim}
  exit
  exit
  \end{verbatim}
\end{enumerate}

\subsubsection{Verifying the Backend}

With the containers running, confirm the API is healthy:

\begin{verbatim}
curl http://localhost:8080/health
\end{verbatim}

A successful response returns \texttt{200 OK} with a JSON body indicating the server and database status.

\subsection{Android Client}

\subsubsection{Using Android Studio (Recommended)}

\begin{enumerate}
  \item Open Android Studio and select \emph{Open an existing project}. Navigate to the \texttt{client/} folder and open it.
  \item Android Studio will detect the Gradle build file and sync dependencies automatically. Accept any SDK or Gradle plugin update prompts.
  \item Create or start an Android Virtual Device (AVD) from the \emph{Device Manager} panel. A device with API level 24 or higher is required (\texttt{minSdk = 24}).
  \item Press \emph{Run} (\texttt{Shift+F10}) to build and install the application on the emulator.
\end{enumerate}

When testing with the emulator, the application connects to the backend at \texttt{http://10.0.2.2:8080}. The special address \texttt{10.0.2.2} is the Android emulator's alias for the host machine's \texttt{localhost}.

\subsubsection{Building from the Command Line}

If Android Studio is unavailable or you prefer the command line:

\begin{enumerate}
  \item Locate your Android SDK directory. On Linux it is typically \texttt{\texttildelow/Android/Sdk}.
  \item Inside \texttt{client/}, create \texttt{local.properties} with one line:
  \begin{verbatim}
  sdk.dir=/home/<your-username>/Android/Sdk
  \end{verbatim}
  \item Initialize Gradle (first time only):
  \begin{verbatim}
  cd client
  ./gradlew
  \end{verbatim}
  \item Build a debug APK:
  \begin{verbatim}
  ./gradlew assembleDebug
  \end{verbatim}
  The APK is written to \texttt{client/app/build/outputs/apk/debug/}.
  \item Install on a connected device or running emulator:
  \begin{verbatim}
  ./gradlew installDebug
  \end{verbatim}
  \item When finished, stop the Gradle daemon to free memory:
  \begin{verbatim}
  ./gradlew --stop
  \end{verbatim}
\end{enumerate}

\subsection{Compiling the Report}

The report is typeset with \LaTeX{} and compiled via \texttt{latexmk}.

\begin{enumerate}
  \item Install \textbf{MikTeX} (Windows/macOS/Linux) or \textbf{TeX Live} (Linux/macOS) and \textbf{Perl}.
  \item Within MikTeX Console, install the \texttt{latexmk} package.
  \item From the \texttt{report/} directory:
  \begin{verbatim}
  latexmk -pvc -pdf main.tex
  \end{verbatim}
  The \texttt{-pvc} flag enables continuous preview mode: \texttt{latexmk} recompiles the PDF whenever a \texttt{.tex} file changes. The compiled PDF is written to \texttt{report/main.pdf}.
\end{enumerate}

\subsection{Stopping the Environment}

To stop the running containers:

\begin{verbatim}
docker compose down
\end{verbatim}

To also remove the persistent database volume (destructive---deletes all seeded data):

\begin{verbatim}
docker compose down -v
\end{verbatim}
